\section{Proofs for Section~\ref{sec:extensions}}

This appendix provides the algebra for the extensions in Section~\ref{sec:extensions}.

\subsection{Preliminaries}

Recall the baseline objects
\[
P_A=\frac{N_A(N_A-1)}{2},
\qquad
Q_{AB}=N_A N_B,
\]
\[
\Delta m_A:=m_A^1-m_A^0>0,
\qquad
\Lambda_A:=(1+\beta)-\beta\delta(m_A^1+m_A^0).
\]

\subsection{Proof of Proposition \ref{prop:partial_internalization}}

Fix $h\in\{0,1\}$ and $\alpha\in[0,1]$. From \eqref{eq:Pi_alpha_compact}, the decentralized benchmark solves
\[
\max_{b\ge 0}
\left\{
\bar\Pi_{\alpha}(h)+\beta A_{h,\alpha}^P b-\frac{\kappa}{2}b^2
\right\}.
\]
Differentiating with respect to $b$ gives
\[
\frac{\partial \Pi_{\alpha}(h,b)}{\partial b}
=
\beta A_{h,\alpha}^P-\kappa b.
\]
Because
\[
\frac{\partial^2 \Pi_{\alpha}(h,b)}{\partial b^2}=-\kappa<0,
\]
the problem is strictly concave and the first-order condition is sufficient. Setting the derivative equal to zero yields
\[
b_{P,\alpha}^*(h)=\frac{\beta A_{h,\alpha}^P}{\kappa},
\]
which is \eqref{eq:bPalpha}.

Next, compare this with the planner's choice from \eqref{eq:bWstar}:
\begin{align}
b_W^*(h)-b_{P,\alpha}^*(h)
&=
\frac{\beta}{\kappa}
\left[
\xi v_LP_Am_A(h)+v_XQ_{AB}(\eta_0+\eta_1 h)
\right]
\nonumber\\
&\quad
-\frac{\beta}{\kappa}
\left[
\alpha\xi v_LP_Am_A(h)+v_XQ_{AB}(\eta_0+\eta_1 h)
\right]
\nonumber\\
&=
\frac{\beta(1-\alpha)\xi v_LP_Am_A(h)}{\kappa},
\label{eq:app7_gap_alpha_raw}
\end{align}
which proves \eqref{eq:b_gap_alpha}. The inequality is strict for every $\alpha<1$ because $\beta,\xi,v_L,P_A,$ and $m_A(h)$ are strictly positive.

Finally, differentiating \eqref{eq:bPalpha} with respect to $\alpha$ gives
\[
\frac{\partial b_{P,\alpha}^*(h)}{\partial \alpha}
=
\frac{\beta\xi v_LP_Am_A(h)}{\kappa}
>0,
\]
which proves \eqref{eq:dbdalpha}. This completes the proof of Proposition \ref{prop:partial_internalization}.
\qed

\subsection{Proof of Proposition \ref{prop:alpha_threshold_gap}}

Recall from Section~\ref{sec:short_run} that
\[
S_A=v_LP_A(m_A^1-m_A^0),
\qquad
D_A=\beta v_LP_A(m_A^1-m_A^0)\bigl[1-\delta(m_A^1+m_A^0)\bigr].
\]
Under partial internalization, the decentralized support threshold is
\[
F_{P,\alpha}^{\dagger}
=
S_A+\alpha D_A+\bigl(R_{1,\alpha}^P-R_{0,\alpha}^P\bigr),
\]
while the planner's threshold remains
\[
F_{H+b}^{\dagger}
=
S_A+D_A+\bigl(R_1-R_0\bigr).
\]
Subtracting the two thresholds gives
\begin{align}
\Omega_{\alpha}
&=
F_{H+b}^{\dagger}-F_{P,\alpha}^{\dagger}
\nonumber\\
&=
(1-\alpha)D_A
\Bigl[(R_1-R_0)-\bigl(R_{1,\alpha}^P-R_{0,\alpha}^P\bigr)\Bigr].
\label{eq:app7_Omega_start}
\end{align}

Write
\[
a_h:=\xi v_LP_Am_h,
\qquad
c_h:=v_XQ_{AB}(\eta_0+\eta_1 h),
\qquad
m_h=m_A^0+(m_A^1-m_A^0)h.
\]
Then
\[
R_h=\frac{\beta^2(a_h+c_h)^2}{2\kappa},
\qquad
R_{h,\alpha}^P=\frac{\beta^2(\alpha a_h+c_h)^2}{2\kappa}.
\]
Hence
\begin{align}
R_h-R_{h,\alpha}^P
&=
\frac{\beta^2}{2\kappa}
\Bigl[(a_h+c_h)^2-(\alpha a_h+c_h)^2\Bigr]
\nonumber\\
&=
\frac{\beta^2}{2\kappa}(1-\alpha)a_h\bigl[(1+\alpha)a_h+2c_h\bigr].
\label{eq:app7_R_gap_h}
\end{align}
Applying \eqref{eq:app7_R_gap_h} at $h=1$ and $h=0$ yields
\begin{align}
(R_1-R_0)-\bigl(R_{1,\alpha}^P-R_{0,\alpha}^P\bigr)
&=
\frac{\beta^2}{2\kappa}(1-\alpha)
\left[
(1+\alpha)(a_1^2-a_0^2)+2(a_1c_1-a_0c_0)
\right].
\label{eq:app7_R_gap_difference}
\end{align}
Now
\[
a_1^2-a_0^2
=
\xi^2 v_L^2P_A^2\bigl((m_A^1)^2-(m_A^0)^2\bigr)
=
X_A,
\]
and
\[
2(a_1c_1-a_0c_0)
=
2\xi v_LP_A v_XQ_{AB}\Bigl[\eta_0(m_A^1-m_A^0)+\eta_1m_A^1\Bigr]
=
Y_A.
\]
Substituting these identities into \eqref{eq:app7_Omega_start} and \eqref{eq:app7_R_gap_difference} gives
\[
\Omega_{\alpha}
=
(1-\alpha)
\left[
D_A+\frac{\beta^2}{2\kappa}\Bigl((1+\alpha)X_A+Y_A\Bigr)
\right],
\]
which is \eqref{eq:Omega_alpha}.

If $\Omega_{\alpha}>0$, then
\[
F_{P,\alpha}^{\dagger}<F_{H+b}^{\dagger},
\]
so every fixed cost in the interval \eqref{eq:alpha_under_support_interval} yields decentralized rejection and planner support. If $\Omega_{\alpha}<0$, the inequality is reversed, giving the interval \eqref{eq:alpha_over_support_interval}. This proves parts (1) and (2).

Part (3) follows immediately from \eqref{eq:Omega_alpha}: setting $\alpha=1$ gives $\Omega_1=0$. Differentiating \eqref{eq:Omega_alpha} with respect to $\alpha$ gives
\begin{align}
\frac{\partial \Omega_{\alpha}}{\partial \alpha}
&=
-D_A-\frac{\beta^2}{2\kappa}\Bigl(2\alpha X_A+Y_A\Bigr).
\label{eq:app7_dOmega_dalpha}
\end{align}
Because $X_A>0$, $Y_A>0$, and $\alpha\in[0,1)$, the second term on the right-hand side is strictly negative. Therefore, if $D_A\ge 0$, then
\[
\frac{\partial \Omega_{\alpha}}{\partial \alpha}<0
\qquad
\text{for all }\alpha\in[0,1),
\]
so the threshold gap is strictly decreasing in $\alpha$. This completes the proof of Proposition \ref{prop:alpha_threshold_gap}.
\qed

\subsection{Proof of Proposition \ref{prop:heterogeneous_participants}}

In the heterogeneous-participant extension, the expected surplus per realized local interaction is
\begin{equation}
\bar v_L(\sigma)
=
v_L^L(1-\sigma)^2
+
v_L^H\bigl[1-(1-\sigma)^2\bigr]
=
v_L^L+(2\sigma-\sigma^2)(v_L^H-v_L^L),
\label{eq:app7_vbarL_sigma}
\end{equation}
with
\begin{equation}
\bar v_L'(\sigma)=2(1-\sigma)(v_L^H-v_L^L)>0
\qquad
\text{for }\sigma\in[0,1).
\label{eq:app7_vbarL_prime}
\end{equation}
Because $\bar v_L(\sigma)$ is a convex combination of $v_L^L$ and $v_L^H$,
\begin{equation}
\bar v_L(\sigma)>0
\qquad
\text{for all }\sigma\in[0,1].
\label{eq:app7_vbarL_positive}
\end{equation}

Likewise,
\begin{equation}
\bar v_X(\sigma)
=
v_X^L(1-\sigma)+v_X^H\sigma
=
v_X^L+\sigma(v_X^H-v_X^L),
\label{eq:app7_vbarX_sigma}
\end{equation}
with
\begin{equation}
\bar v_X'(\sigma)=v_X^H-v_X^L>0
\label{eq:app7_vbarX_prime}
\end{equation}
and
\begin{equation}
\bar v_X(\sigma)>0
\qquad
\text{for all }\sigma\in[0,1].
\label{eq:app7_vbarX_positive}
\end{equation}

Using \eqref{eq:vbarL_sigma}, the hub-alone welfare gain is
\[
\Delta_H(\sigma)=\bar v_L(\sigma)P_A\Delta m_A\Lambda_A-F.
\]
Differentiating with respect to $\sigma$ gives
\begin{align}
\frac{\partial \Delta_H(\sigma)}{\partial \sigma}
&=
\bar v_L'(\sigma)P_A\Delta m_A\Lambda_A.
\label{eq:app7_dDeltaH_dsigma_raw}
\end{align}
If $\Lambda_A>0$, then \eqref{eq:app7_vbarL_prime} implies
\[
\frac{\partial \Delta_H(\sigma)}{\partial \sigma}>0
\qquad
\text{for }\sigma\in[0,1),
\]
which proves \eqref{eq:dDeltaH_dsigma}.

Next, from \eqref{eq:bstar_sigma},
\[
b^*(\sigma)=\frac{\beta A_1(\sigma)}{\kappa},
\qquad
A_1(\sigma)=\xi \bar v_L(\sigma)P_A m_A^1+\bar v_X(\sigma)Q_{AB}(\eta_0+\eta_1).
\]
Differentiating,
\begin{align}
\frac{\partial b^*(\sigma)}{\partial \sigma}
&=
\frac{\beta}{\kappa}
\left[
\xi \bar v_L'(\sigma)P_A m_A^1
+
\bar v_X'(\sigma)Q_{AB}(\eta_0+\eta_1)
\right].
\label{eq:app7_dbstar_dsigma_raw}
\end{align}
Since all multiplicative terms are positive, \eqref{eq:app7_vbarL_prime} and \eqref{eq:app7_vbarX_prime} imply
\[
\frac{\partial b^*(\sigma)}{\partial \sigma}>0,
\]
which proves \eqref{eq:dbstar_dsigma}.

For the renewal surplus,
\[
R_1(\sigma)=\frac{\beta^2A_1(\sigma)^2}{2\kappa}.
\]
Differentiating,
\begin{align}
\frac{\partial R_1(\sigma)}{\partial \sigma}
&=
\frac{\beta^2}{\kappa}
A_1(\sigma)
\left[
\xi \bar v_L'(\sigma)P_A m_A^1
+
\bar v_X'(\sigma)Q_{AB}(\eta_0+\eta_1)
\right].
\label{eq:app7_dR_dsigma_raw}
\end{align}
By \eqref{eq:app7_vbarL_positive}, \eqref{eq:app7_vbarX_positive}, \eqref{eq:app7_vbarL_prime}, and \eqref{eq:app7_vbarX_prime}, every factor on the right-hand side is strictly positive. Hence
\[
\frac{\partial R_1(\sigma)}{\partial \sigma}>0,
\]
which proves \eqref{eq:dR_dsigma}. This completes the proof of Proposition \ref{prop:heterogeneous_participants}.
\qed
